\chapter{光滑流形}

\begin{solution}{manifold-sphere-stereographic}
记 \(N=(0,\ldots,0,1)\)、\(S=(0,\ldots,0,-1)\)，并把球面点写成
\((u,t)\in\R^n\times\R\)。取
\[
 \sigma_N(u,t)=\frac{u}{1-t},\qquad
 \sigma_S(u,t)=\frac{u}{1+t}.
\]
直接解方程得到
\[
 \sigma_N^{-1}(x)=
 \left(\frac{2x}{1+\norm{x}^2},
       \frac{\norm{x}^2-1}{1+\norm{x}^2}\right),
 \qquad
 \sigma_S^{-1}(y)=
 \left(\frac{2y}{1+\norm{y}^2},
       \frac{1-\norm{y}^2}{1+\norm{y}^2}\right).
\]
两个分母都恒正；代入后可以直接验证所得点在 \(S^n\) 上，并分别避开
\(N,S\)。两图重叠对应 \(x\ne0\)，而
\[
 (\sigma_S\circ\sigma_N^{-1})(x)=\frac{x}{\norm{x}^2}.
\]
这个映射在 \(\R^n\setminus\{0\}\) 上光滑，并且等于自己的逆映射，故
两张立体投影图光滑相容。
\end{solution}

\begin{solution}{manifold-product-atlas}
若 \((U_i,x_i)\)、\((V_j,y_j)\) 分别来自两个光滑图册，则两张积图之间
的坐标变换为
\[
 (x_k,y_\ell)\circ(x_i,y_j)^{-1}(u,v)
 =\bigl(x_k\circ x_i^{-1}(u),\,
        y_\ell\circ y_j^{-1}(v)\bigr)
\]
（等式限于相应公共定义域）。两个分量都光滑，逆变换同理，所以积图构成
光滑图册。

映射
\[
 ((p,q),r)\longmapsto(p,(q,r))
\]
是拓扑同胚，在任意积坐标中只改变括号而不改变欧氏坐标，因而其坐标表示
是恒等线性映射。它和逆映射都光滑，给出
\[
 (M\times N)\times P\cong M\times(N\times P)
\]
之间的自然微分同胚。
\end{solution}

\begin{solution}{manifold-atlas-common-refinement}
若图册 \(\mathcal A,\mathcal B\) 等价，对每个
\((U,x)\in\mathcal A\)、\((V,y)\in\mathcal B\) 的非空交集
\(U\cap V\)，取限制图
\[
 (U\cap V,x|_{U\cap V}).
\]
这些图覆盖 \(M\)，其坐标域同时细于两组原坐标域。它们彼此光滑相容，并且
同 \(\mathcal A,\mathcal B\) 中所有图相容，故构成共同光滑加细。

反过来，设 \(\mathcal C\) 是这样的共同光滑加细。取
\((U,x)\in\mathcal A\)、\((V,y)\in\mathcal B\) 及
\(p\in U\cap V\)。从 \(\mathcal C\) 中取含 \(p\) 的图 \((W,z)\)，
并把 \(W\) 缩小到 \(U\cap V\) 内。在 \(p\) 的一个邻域上，
\[
 y\circ x^{-1}
 =(y\circ z^{-1})\circ(z\circ x^{-1}).
\]
右端两个因子光滑。坐标变换的光滑性是局部性质，所以 \(x,y\) 在整个
重叠处相容。任取两图均如此，故两个图册等价。
\end{solution}

\begin{solution}{manifold-smoothness-all-charts}
设 \(F\) 在给定的源、目标覆盖图中都有光滑坐标表示。任取极大图册中的
源图 \((U,x)\) 和目标图 \((V,y)\)。对每个
\[
 p\in U\cap F^{-1}(V),
\]
从给定覆盖中取图 \((U_0,x_0)\)、\((V_0,y_0)\)，分别覆盖
\(p,F(p)\)。缩小 \(p\) 的邻域，使其像落在 \(V\cap V_0\) 中，则
\[
 y\circ F\circ x^{-1}
 =(y\circ y_0^{-1})\circ
  (y_0\circ F\circ x_0^{-1})\circ
  (x_0\circ x^{-1}).
\]
三个因子都光滑，故左侧在每一点附近光滑，从而在其定义域上光滑。
\end{solution}

\begin{solution}{manifold-rank-coordinate-invariance}
若 \(\widetilde x,\widetilde y\) 是另一对源、目标坐标，则链式法则给出
\[
 D(\widetilde y\circ F\circ\widetilde x^{-1})
 =D(\widetilde y\circ y^{-1})\,
  D(y\circ F\circ x^{-1})\,
  D(x\circ\widetilde x^{-1}).
\]
左右两个坐标变换的导数都可逆。矩阵左乘或右乘可逆矩阵不改变秩，因此
中间矩阵的秩与坐标选择无关。浸入和浸没的定义也就与坐标无关。
\end{solution}

\begin{solution}{manifold-graph-submanifold}
定义
\[
 \Phi:U\times\R^{n-k}\longrightarrow U\times\R^{n-k},\qquad
 \Phi(u,v)=(u,v-f(u)).
\]
它的逆映射为
\[
 \Phi^{-1}(u,w)=(u,w+f(u)),
\]
所以 \(\Phi\) 是微分同胚，并且
\[
 \Phi(\Gamma_f)=U\times\{0\}.
\]
这给出切片坐标，证明 \(\Gamma_f\) 是嵌入子流形。

投影
\[
 \pi:\Gamma_f\longrightarrow U,\qquad (u,f(u))\longmapsto u
\]
光滑，其逆映射 \(u\mapsto(u,f(u))\) 也光滑，故 \(\pi\) 是微分同胚。
\end{solution}

\begin{solution}{manifold-sphere-regular-value}
令 \(F(x)=\norm{x}^2\)。对 \(x\in S^{n-1}\)，
\[
 DF(x)h=2x\cdot h.
\]
取 \(h=x\) 得 \(DF(x)x=2\)，故线性映射
\(DF(x):\R^n\to\R\) 满射。于是 \(1\) 是正则值，正则值定理说明
\(S^{n-1}\) 是 \(n-1\) 维嵌入子流形。

给定 \(p\in S^{n-1}\)，可选 \(j\) 使 \(p_j\ne0\)。在 \(p\) 的
足够小邻域中固定第 \(j\) 个坐标的符号，球面可写成
\[
 x_j=\operatorname{sgn}(p_j)
 \sqrt{1-\sum_{i\ne j}x_i^2}.
\]
根号内在这个邻域中为正，其余 \(n-1\) 个坐标便给出所求局部图像坐标。
\end{solution}

\begin{solution}{manifold-special-linear-submanifold}
把矩阵空间识别为 \(\R^{n^2}\)，令 \(F(A)=\det A\)。当 \(A\) 可逆时，
Jacobi 公式给出
\[
 DF(A)H=\det(A)\operatorname{tr}(A^{-1}H).
\]
若 \(A\in SL(n,\R)\)，取 \(H=A\)，便有
\[
 DF(A)A=n\ne0.
\]
因此 \(DF(A):\R^{n^2}\to\R\) 满射。正则值定理说明
\[
 F^{-1}(1)=SL(n,\R)
\]
是 \(\R^{n^2}\) 中的 \(n^2-1\) 维嵌入子流形。
\end{solution}

\begin{solution}{manifold-singular-level-examples}
方程 \(x^2-y^2=0\) 等价于 \(y=x\) 或 \(y=-x\)。零水平集在原点是两条
直线交叉的并。去掉原点后，原点的任意充分小邻域与该集合的交有四个连通
分支；而开区间去掉一点只有两个连通分支。因此该集合在原点附近不同胚于
开区间，不是一维流形。梯度
\[
 \nabla(x^2-y^2)=(2x,-2y)
\]
在原点为零。

对 \(F(x)=x^3\)，零水平集只是 \(\{0\}\)，它是 \(\R\) 中的零维嵌入
子流形；但 \(F'(0)=0\)，所以 \(0\) 不是正则值。这两个例子分别说明：
正则值条件失效时，水平集可能出现奇点，也可能仍然是子流形。
\end{solution}

\begin{solution}{manifold-injective-immersion-not-embedding}
首先
\[
 \gamma'(t)=(\cos t,2\cos 2t).
\]
若第一分量为零，则 \(t=\pm\pi/2\)，此时第二分量等于 \(-2\)。所以
\(\gamma'(t)\) 从不为零，\(\gamma\) 是浸入。

再证单射。若 \(\sin s=\sin t\)，在 \((-\pi,\pi)\) 中除 \(s=t\) 外，
只可能有 \(s=\pi-t\) 或 \(s=-\pi-t\)，只取其中仍落在给定区间的一项。
在这两种情形下都有 \(\sin 2s=-\sin 2t\)。若第二坐标也相等，则
\(\sin 2t=0\)。检查 \(t=-\pi/2,0,\pi/2\) 后，或者另一个候选点落在
区间外，或者它就是 \(t\) 本身。因此 \(s=t\)，从而 \(\gamma\) 单射。

然而，令 \(t_k=\pi-1/k\)。当 \(k\) 足够大时 \(t_k\in(-\pi,\pi)\)，且
\[
 \gamma(t_k)\longrightarrow(0,0)=\gamma(0),
 \qquad t_k\not\longrightarrow0.
\]
若 \(\gamma^{-1}:\gamma((-\pi,\pi))\to(-\pi,\pi)\) 连续，上述收敛将
迫使 \(t_k\to0\)，矛盾。因此 \(\gamma\) 不是到其像的同胚，也就不是
光滑嵌入。
\end{solution}

\begin{solution}{manifold-cutoff-finite-combination}
设 \(0\le\chi_j\le1\)，每个 \(\chi_j\) 在覆盖 \(K\) 的某个开集上
等于 \(1\)，并且 \(\operatorname{supp}\chi_j\Subset U\)。有限乘积
\[
 \chi=1-\prod_{j=1}^r(1-\chi_j)
\]
光滑且取值于 \([0,1]\)。若一点属于某个 \(\chi_j=1\) 的邻域，则
\(\chi=1\)。这些邻域覆盖 \(K\)，其并是 \(K\) 的开邻域，所以
\(\chi=1\) 于 \(K\) 的一个邻域。

在所有 \(\operatorname{supp}\chi_j\) 之外，各 \(\chi_j\) 都为零，
因而
\[
 \operatorname{supp}\chi
 \subset\bigcup_{j=1}^r\operatorname{supp}\chi_j\Subset U.
\]
右端是有限个紧集之并。每个局部截断都已经在坐标域外光滑补零，有限次乘法
和减法保持光滑，故 \(\chi\in C_c^\infty(M)\)。
\end{solution}

\begin{solution}{manifold-locally-finite-smooth-sum}
若 \(\{f_i\}\) 局部有限，则每个 \(p\) 都有邻域 \(V\)，其上只有有限多个
\(f_i\) 非零。因此
\[
 \left.\sum_i f_i\right|_V
\]
是有限个光滑函数之和，故光滑；这些局部表达在重叠处显然一致。

反例可在 \(\R\) 上构造。对 \(n\ge2\)，令
\[
 I_n=
 \left(
 \frac12\left(\frac1{n+1}+\frac1n\right),
 \frac12\left(\frac1n+\frac1{n-1}\right)
 \right).
\]
这些开区间两两不交，并且 \(1/n\in I_n\)。取
\(\varphi_n\in C_c^\infty(I_n)\) 满足 \(\varphi_n(1/n)=1\)。每一点
至多属于一个支集，所以逐点只有有限项非零；但这族在 \(0\) 的任何邻域中
都有无穷多个非零成员，因而不局部有限。逐点和
\[
 f=\sum_{n\ge2}\varphi_n
\]
满足 \(f(0)=0\)、\(f(1/n)=1\)，故在 \(0\) 不连续。
\end{solution}

\begin{solution}{manifold-partition-glue-functions}
先把 \(\rho_i f_i\) 在 \(U_i\) 外定义为零。由于
\(\operatorname{supp}\rho_i\subset U_i\)，每一点若在 \(U_i\) 的边界
附近，则有一个邻域不遇到 \(\operatorname{supp}\rho_i\)；所以补零函数
仍然光滑。单位分解局部有限，因此
\[
 f=\sum_i\rho_i f_i
\]
在每一点附近都是有限和，故良定义且光滑。

它一般只是局部函数的加权平均。若两个重叠开集上分别取常函数
\(f_1=0\)、\(f_2=1\)，并且某点满足 \(0<\rho_2<1\)，则该点有
\(f=\rho_2\)，既不等于 \(f_1\) 也不等于 \(f_2\)。只有各局部函数在
重叠处相等时，单位分解粘合才恢复它们共同定义的函数。
\end{solution}

\begin{solution}{manifold-compact-embedding-reconstruction}
取有限坐标覆盖 \(\{(U_i,x_i)\}_{i=1}^m\)，以及从属于它的单位分解
\(\{\rho_i\}_{i=1}^m\)，定义
\[
 \Phi(p)=
 \bigl(\rho_1(p),\rho_1(p)x_1(p),\ldots,
       \rho_m(p),\rho_m(p)x_m(p)\bigr).
\]
各 \(\rho_i x_i\) 因 \(\operatorname{supp}\rho_i\Subset U_i\) 而可在
\(U_i\) 外光滑补零。

若 \(\Phi(p)=\Phi(q)\)，选 \(i\) 使 \(\rho_i(p)>0\)。标量分量给出
\(\rho_i(q)=\rho_i(p)>0\)，所以 \(p,q\in U_i\)；向量分量随后给出
\(x_i(p)=x_i(q)\)，故 \(p=q\)。

再在 \(p\) 处用 \(x_i\) 作坐标。若
\[
 D(\Phi\circ x_i^{-1})(x_i(p))v=0,
\]
则第 \(i\) 个标量分量给 \(D\rho_i(p)v=0\)，相应向量分量给
\[
 \rho_i(p)v+x_i(p)D\rho_i(p)v=\rho_i(p)v=0.
\]
于是 \(v=0\)，所以 \(\Phi\) 是浸入。最后，\(M\) 紧而欧氏空间
Hausdorff，连续单射 \(\Phi\) 是到其像的同胚。因此 \(\Phi\) 是光滑
嵌入。
\end{solution}
